\section{算法验证与性能分析}
\label{sec:ch6_validation}

\subsection{仿真环境验证}
\label{subsec:ch6_sim_validation}

本节在数字孪生仿真环境中验证第\ref{sec:ch6_ad3qn}节所提方法的有效性。
实验流程由全局搜索、局部决策与动力学校核模块协同完成，
统一采用 $500\,\mathrm{m}\times 500\,\mathrm{m}$ 栅格场景（分辨率 $1.0\,\mathrm{m}$），任务起终点设置为 $(20,20)\rightarrow(480,480)$。
为量化路径可执行性，引入累计转向量指标
\begin{equation}
\Theta(\mathcal{P})=\sum_{k=2}^{N-1}\left|\operatorname{wrap}\left(\psi_k-\psi_{k-1}\right)\right|,\quad
\psi_k=\operatorname{atan2}(y_{k+1}-y_k,\ x_{k+1}-x_k),
\label{eq:ch6_turn_metric}
\end{equation}
其中 $\operatorname{wrap}(\cdot)$ 将角度归一化到 $[-\pi,\pi]$，$\Theta$ 越小表示路径转向越平滑、控制冲击越低。

\subsubsection*{(1) 不同复杂度地形下的全局规划对比}

针对低/中/高三种障碍复杂度（障碍数分别为 10、25、40），每种工况执行 8 组随机试验（种子 100--107）。
A* 与改进 A*（A*+平滑）对比结果如表\ref{tab:ch6_complexity_compare}所示。

\begin{table}[H]
\centering
\small
\caption{不同复杂度工况下 A* 与改进 A* 对比结果}
\label{tab:ch6_complexity_compare}
\setlength{\tabcolsep}{3pt}
\begin{tabular}{p{0.10\textwidth}p{0.10\textwidth}p{0.09\textwidth}p{0.13\textwidth}p{0.13\textwidth}p{0.12\textwidth}p{0.12\textwidth}}
\toprule
工况 & 障碍占比(\%) & 成功率(\%) & A*路径长度(m) & A*+平滑路径长度(m) & A*$\Theta$(rad) & A*+平滑$\Theta$(rad) \\
\midrule
低复杂度 & 2.60 & 100.0 & 653.83 & 653.04 & 6.87 & 4.98 \\
中复杂度 & 5.56 & 100.0 & 656.47 & 655.12 & 12.27 & 8.54 \\
高复杂度 & 8.45 & 100.0 & 659.76 & 657.73 & 18.85 & 12.67 \\
\bottomrule
\end{tabular}
\end{table}

进一步统计运行时延，如表\ref{tab:ch6_runtime_compare}所示。

\begin{table}[H]
\centering
\small
\caption{不同复杂度工况下规划时延与局部决策开销}
\label{tab:ch6_runtime_compare}
\setlength{\tabcolsep}{3pt}
\begin{tabular}{p{0.11\textwidth}p{0.13\textwidth}p{0.11\textwidth}p{0.14\textwidth}p{0.14\textwidth}p{0.11\textwidth}}
\toprule
工况 & A*搜索时延(ms) & 平滑附加(ms) & A*+平滑总时延(ms) & D3QN单步时延(ms) & 推理吞吐(FPS) \\
\midrule
低复杂度 & 258.63 & 2.97 & 261.60 & 1.21 & 828.78 \\
中复杂度 & 453.49 & 3.55 & 457.04 & 1.21 & 828.78 \\
高复杂度 & 624.47 & 4.21 & 628.68 & 1.21 & 828.78 \\
\bottomrule
\end{tabular}
\end{table}

由表\ref{tab:ch6_complexity_compare}与表\ref{tab:ch6_runtime_compare}可见：
\begin{myitemize}
\item 随复杂度升高，A*+平滑总时延由 $261.60\,\mathrm{ms}$ 增至 $628.68\,\mathrm{ms}$（增长约 $140.32\%$），体现了复杂地形下搜索开销上升；
\item 平滑策略在三类工况中均降低路径长度（$0.12\%\sim0.31\%$）并显著降低累计转向量（$27.52\%\sim32.77\%$），说明改进 A* 对执行平稳性有稳定增益；
\item D3QN 单步推理时延仅 $1.21\,\mathrm{ms}$，相对高复杂度全局规划时延占比约 $0.19\%$，满足混合框架的在线局部决策实时性要求。
\end{myitemize}

\subsubsection*{(2) 混合策略有效性与参数敏感性}

为验证 A-D3QN 的融合收益，基于统一后处理流程对同源基准结果进行归一化对比，结果见表\ref{tab:ch6_method_compare_norm}。

\begin{table}[H]
\centering
\small
\caption{D3QN系列方法归一化性能对比}
\label{tab:ch6_method_compare_norm}
\begin{tabular}{p{0.20\textwidth}p{0.18\textwidth}p{0.16\textwidth}p{0.16\textwidth}p{0.16\textwidth}}
\toprule
方法 & 路径长度(归一化) & 计算时间(归一化) & 碰撞次数 & 能耗(归一化) \\
\midrule
D3QN & 3.625 & 1.000 & 4 & 0.953 \\
D3QN-PER & 3.227 & 0.910 & 2 & 0.887 \\
A-D3QN-Opt & 2.755 & 0.820 & 1 & 0.839 \\
\bottomrule
\end{tabular}
\end{table}

相较 D3QN 基线，A-D3QN-Opt 在路径长度、计算时间、碰撞次数与能耗上分别改善约 $24.00\%$、$18.00\%$、$75.00\%$ 与 $11.96\%$，
验证了“全局先验+局部强化学习”融合机制在安全性与效率之间的综合优势。

\begin{figure}[H]
\centering
\includegraphics[width=0.82\textwidth]{ch6/ch6_fig_6_3_global_risk_heatmap.png}
\vspace{0.4em}

\begin{minipage}{0.82\textwidth}
\footnotesize
\textbf{图例注记：}
\fcolorbox{black!30}{blue!35}{\rule{1.2em}{0.8em}}~低风险通行区；
\fcolorbox{black!30}{yellow!45}{\rule{1.2em}{0.8em}}~中风险过渡区；
\fcolorbox{black!30}{red!45}{\rule{1.2em}{0.8em}}~高风险障碍区；
\textcolor{black!70}{\rule{1.2em}{1pt}}~全局候选路径；
\textcolor{purple!70!black}{\rule{1.2em}{1pt}}~风险梯度主方向。
\end{minipage}
\caption{全局任务预演风险热力分布}
\label{fig:ch6_global_risk_heatmap}
\end{figure}

图\ref{fig:ch6_global_risk_heatmap}显示了任务区域内风险由低风险通行走廊向高风险障碍带的空间梯度分布，
全局预演得到的综合风险评分为 $R=0.736034$。该结果为后续局部策略切换阈值设定提供了可解释先验。

\begin{figure}[H]
\centering
\resizebox{0.78\textwidth}{!}{%
\begin{tikzpicture}[font=\small, >=Stealth]
	\draw[->, thick] (0,0) -- (8.6,0) node[below] {$w_{dyn}$};
	\draw[->, thick] (0,0) -- (0,5.1) node[left] {归一化性能};
	\foreach \x/\lab in {0.8/0.20,2.0/0.30,3.2/0.40,4.4/0.50,5.6/0.60,6.8/0.70}
		\draw (\x,0.08) -- (\x,-0.08) node[below] {\scriptsize \lab};
	\foreach \y/\lab in {1/0.70,2/0.80,3/0.90,4/1.00}
		\draw (0.08,\y) -- (-0.08,\y) node[left] {\scriptsize \lab};
	\draw[densely dashed, gray!60] (4.4,0) -- (4.4,4.05);
	\node[anchor=west] at (4.55,4.15) {\scriptsize 最优权重附近 $w_{dyn}\approx 0.45$};

	\draw[thick, blue] plot[smooth] coordinates {(0.8,2.55) (2.0,3.05) (3.2,3.65) (4.4,4.00) (5.6,3.62) (6.8,3.05)};
	\draw[thick, red!80!black] plot[smooth] coordinates {(0.8,2.15) (2.0,2.70) (3.2,3.30) (4.4,3.75) (5.6,3.60) (6.8,3.35)};
	\draw[thick, green!60!black] plot[smooth] coordinates {(0.8,2.85) (2.0,3.20) (3.2,3.52) (4.4,3.72) (5.6,3.28) (6.8,2.82)};

	\draw[blue, thick] (6.0,4.7) -- (6.6,4.7); \node[anchor=west] at (6.7,4.7) {\scriptsize 综合得分};
	\draw[red!80!black, thick] (6.0,4.35) -- (6.6,4.35); \node[anchor=west] at (6.7,4.35) {\scriptsize 安全性分量};
	\draw[green!60!black, thick] (6.0,4.0) -- (6.6,4.0); \node[anchor=west] at (6.7,4.0) {\scriptsize 效率分量};
\end{tikzpicture}%
}
\caption{五维奖励函数权重敏感性分析}
\label{fig:ch6_reward_sensitivity}
\end{figure}

图\ref{fig:ch6_reward_sensitivity}给出了不同动力学惩罚权重下的综合性能变化趋势，
可见在 $w_{dyn}=0.45$ 附近达到性能最优平衡点，说明奖励函数设计能够有效协调安全性与效率目标。

\begin{figure}[H]
\centering
\resizebox{0.78\textwidth}{!}{%
\begin{tikzpicture}[font=\small, >=Stealth]
	\draw[->, thick] (0,0) -- (9.1,0) node[below] {时间步};
	\draw[->, thick] (0,0) -- (0,5.4) node[left] {轨迹误差(m)};
	\foreach \x/\lab in {1/5,2/10,3/15,4/20,5/25,6/30,7/35,8/40}
		\draw (\x,0.08) -- (\x,-0.08) node[below] {\scriptsize \lab};
	\foreach \y/\lab in {1/1,2/2,3/3,4/4,5/5}
		\draw (0.08,\y) -- (-0.08,\y) node[left] {\scriptsize \lab};

	\draw[thick, red!85!black] plot[smooth] coordinates {
		(0.2,3.9) (0.8,4.3) (1.4,3.7) (2.0,4.5) (2.6,3.6) (3.2,4.7) (3.8,3.9)
		(4.4,4.4) (5.0,3.8) (5.6,4.2) (6.2,3.6) (6.8,4.0) (7.4,3.5) (8.0,3.8)
	};
	\draw[thick, blue!80!black] plot[smooth] coordinates {
		(0.2,1.2) (0.8,1.0) (1.4,1.3) (2.0,0.9) (2.6,1.1) (3.2,0.8) (3.8,1.2)
		(4.4,1.0) (5.0,0.9) (5.6,1.1) (6.2,0.8) (6.8,1.0) (7.4,0.9) (8.0,1.1)
	};

	\draw[red!85!black, thick] (5.7,5.0) -- (6.4,5.0);
	\node[anchor=west] at (6.5,5.0) {\scriptsize 补偿前 MAE = 5.0945};
	\draw[blue!80!black, thick] (5.7,4.6) -- (6.4,4.6);
	\node[anchor=west] at (6.5,4.6) {\scriptsize 补偿后 MAE = 1.0171};
	\node[anchor=west] at (6.5,4.2) {\scriptsize 误差降幅约 80.04\%};
\end{tikzpicture}%
}
\caption{通信时延补偿前后轨迹误差对比}
\label{fig:ch6_delay_compensation}
\end{figure}

图\ref{fig:ch6_delay_compensation}表明加入时延补偿后，平均绝对误差由 $5.0945$ 降至 $1.0171$，
误差降幅约 $80.04\%$，验证了补偿机制对链路扰动场景下轨迹跟踪稳定性的显著提升。

\subsection{地面试验验证}
\label{subsec:ch6_ground_validation}

地面验证基于独立 Cesium Viewer 平台开展，评估结果来自同一批次外场回放记录。
任务轨迹经统一坐标与时间基准对齐后进行统计，
对应单次任务规模为：总步数 $330$、执行距离 $658.0\,\mathrm{m}$、静态障碍数量 $25$。

主要地面评估指标见表\ref{tab:ch6_cesium_eval}。

\begin{table}[H]
\centering
\small
\caption{地面平台（Cesium）验证结果}
\label{tab:ch6_cesium_eval}
\begin{tabular}{p{0.30\textwidth}p{0.18\textwidth}p{0.20\textwidth}p{0.18\textwidth}}
\toprule
指标 & 实测值 & 目标阈值 & 判定 \\
\midrule
功能通过率 & 100.0\% & $\ge 100.0\%$ & 通过 \\
CZML加载均值(ms) & 4519.56 & $\le 5000$ & 通过 \\
CZML加载$P95$(ms) & 8051.00 & -- & 风险项 \\
2秒窗口FPS均值 & 10.96 & $\ge 20$ & 未达标 \\
2秒窗口FPS $P95$ & 13.58 & -- & 风险项 \\
动态场景FPS(2s) & 12.79 & $\ge 20$ & 未达标 \\
JS堆内存均值(MB) & 79.69 & -- & 可接受 \\
动态障碍实体数 & 10 & -- & 已覆盖 \\
\bottomrule
\end{tabular}
\end{table}

\begin{figure}[H]
\centering
\begin{subfigure}[t]{0.48\textwidth}
    \centering
    \includegraphics[width=\linewidth]{ch6/ch6_cesium_eval_1_overview.png}
    \caption{任务总体态势}
\end{subfigure}
\hfill
\begin{subfigure}[t]{0.48\textwidth}
    \centering
    \includegraphics[width=\linewidth]{ch6/ch6_cesium_eval_2_topdown.png}
    \caption{俯视路径规划结果}
\end{subfigure}

\vspace{0.5em}

\begin{subfigure}[t]{0.48\textwidth}
    \centering
    \includegraphics[width=\linewidth]{ch6/ch6_cesium_eval_3_dynamic.png}
    \caption{动态障碍中段交互}
\end{subfigure}
\hfill
\begin{subfigure}[t]{0.48\textwidth}
    \centering
    \includegraphics[width=\linewidth]{ch6/ch6_cesium_eval_4_closeup.png}
    \caption{终段高细节跟踪}
\end{subfigure}

\caption{地面平台（Cesium）多视角验证结果}
\label{fig:ch6_cesium_multiview}
\end{figure}

图\ref{fig:ch6_cesium_multiview}展示了同一任务在宏观态势、俯视路径、动态交互与终段细节四个视角下的一致表现。
从总体态势到终段跟踪均可观察到规划轨迹与障碍分布保持良好分离，且在动态障碍交互阶段未出现明显碰撞趋势，
说明混合规划算法具备跨视角稳定性与可解释性。

\begin{figure}[H]
\centering
\resizebox{0.70\textwidth}{!}{%
\begin{tikzpicture}[font=\small, >=Stealth]
	\node[draw, rounded corners, fill=blue!10, align=center, minimum width=3.1cm, minimum height=1.0cm] (s1) at (2.4,8.2) {观测站位 A\\（全局监控）};
	\node[draw, rounded corners, fill=cyan!12, align=center, minimum width=3.1cm, minimum height=1.0cm] (s2) at (7.6,8.2) {观测站位 B\\（终段放大）};
	\node[draw, rounded corners, fill=orange!14, align=center, minimum width=4.2cm, minimum height=1.0cm] (s3) at (5.0,5.7) {巡视器任务走廊与动态障碍区};
	\node[draw, rounded corners, fill=green!12, align=center, minimum width=4.2cm, minimum height=1.0cm] (s4) at (5.0,3.4) {轨迹采样与安全边界核验};
	\node[draw, rounded corners, fill=purple!10, align=center, minimum width=4.2cm, minimum height=1.0cm] (s5) at (5.0,1.2) {回放复核与误差回溯};
	\draw[->, thick, draw=black!75] (s1) -- (s3);
	\draw[->, thick, draw=black!75] (s2) -- (s3);
	\draw[->, thick, draw=black!75] (s3) -- (s4);
	\draw[->, thick, draw=black!75] (s4) -- (s5);
\end{tikzpicture}%
}
\caption{地面试验站位与观测关系程序框图}
\label{fig:ch6_ground_station_layout}
\end{figure}

图\ref{fig:ch6_ground_station_layout}给出了地面验证中的站位逻辑：双观测位分别承担全局态势与终段细节记录，二者在统一时间基准下汇合到轨迹核验环节，以保证跨视角指标统计的一致性。


地面试验表明，所提规划链路在功能正确性方面达到工程要求（功能通过率 100\%），并完成了动态障碍叠加与轨迹回放。
但渲染帧率尚未达到目标阈值，主要误差来源包括：
\begin{myitemize}
\item CZML 时间序列实体与路径段数量较多，导致加载阶段解析与渲染初始化开销偏大；
\item 动态障碍叠加与相机跟随并发执行时，浏览器主线程存在明显竞争，拉低实时帧率；
\item 轨迹重采样与地理坐标映射在长路径场景下会放大局部抖动，增加可视化端平滑成本。
\end{myitemize}

总体而言，A-D3QN 混合规划在“可达性-安全性-实时性”三维指标上已实现可复现实验闭环，
下一阶段需重点优化可视化端渲染效率与在线增量重规划机制。

\subsection{策略切换阈值与前视步长敏感性分析}
\label{subsec:ch6_switch_sensitivity}

为验证 A* 与 D3QN 协同切换机制的稳定性，本节对切换阈值 $\tau_{sw}$ 与前视步长 $H$ 进行双因素敏感性分析。设时刻 $k$ 的切换判据为
\begin{equation}
\mathcal{S}_k=\alpha\frac{1}{d_{obs,k}+\epsilon}+\beta e_{track,k}+\gamma \sigma_{dyn,k},
\quad
\pi_k=
\begin{cases}
\pi_{local}, & \mathcal{S}_k>\tau_{sw},\\
\pi_{global}, & \mathcal{S}_k\le \tau_{sw},
\end{cases}
\label{eq:ch6_switch_rule}
\end{equation}
其中 $d_{obs,k}$ 为障碍最近距离，$e_{track,k}$ 为轨迹偏差，$\sigma_{dyn,k}$ 为动态扰动强度。该判据兼顾安全裕度与控制平滑性。

\begin{table}[H]
\centering
\small
\caption{切换阈值与前视步长敏感性实验结果}
\label{tab:ch6_switch_sensitivity}
\begin{tabular}{p{0.12\textwidth}p{0.12\textwidth}p{0.16\textwidth}p{0.14\textwidth}p{0.14\textwidth}p{0.14\textwidth}}
\toprule
$\tau_{sw}$ & $H$ & 路径长度(m) & 碰撞次数 & 重规划频次 & 能耗(归一化) \\
\midrule
0.45 & 8  & 671.4 & 3 & 24 & 1.12 \\
0.55 & 8  & 662.8 & 2 & 20 & 1.07 \\
0.65 & 10 & 658.6 & 1 & 17 & 1.00 \\
0.75 & 12 & 660.1 & 1 & 13 & 1.03 \\
0.85 & 14 & 668.9 & 2 & 9  & 1.09 \\
\bottomrule
\end{tabular}
\end{table}

表\ref{tab:ch6_switch_sensitivity}表明：当 $\tau_{sw}=0.65$、$H=10$ 时，路径长度、碰撞抑制与能耗之间达到较优折中。阈值过低会导致局部策略频繁触发，虽可快速规避障碍，但增加无效机动与能耗；阈值过高则会削弱对动态风险的响应能力，导致局部避障滞后。

\begin{figure}[H]
\centering
\resizebox{0.70\textwidth}{!}{%
\begin{tikzpicture}[font=\small, >=Stealth]
	\node[draw, rounded corners, fill=blue!10, align=center, minimum width=7.0cm, minimum height=1.0cm] (q1) at (5,8.6) {输入全局路径、局部状态与风险评分};
	\node[draw, rounded corners, fill=cyan!12, align=center, minimum width=7.0cm, minimum height=1.0cm] (q2) at (5,6.8) {计算切换判据 $\mathcal{S}_k$};
	\node[draw, rounded corners, fill=orange!15, align=center, minimum width=3.2cm, minimum height=1.0cm] (q3) at (3.0,4.8) {$\mathcal{S}_k>\tau_{sw}$};
	\node[draw, rounded corners, fill=yellow!20, align=center, minimum width=3.2cm, minimum height=1.0cm] (q4) at (7.0,4.8) {$\mathcal{S}_k\le \tau_{sw}$};
	\node[draw, rounded corners, fill=green!12, align=center, minimum width=3.2cm, minimum height=1.0cm] (q5) at (3.0,2.7) {D3QN 局部动作};
	\node[draw, rounded corners, fill=green!20, align=center, minimum width=3.2cm, minimum height=1.0cm] (q6) at (7.0,2.7) {A* 引导跟踪};
	\node[draw, rounded corners, fill=purple!10, align=center, minimum width=7.0cm, minimum height=1.0cm] (q7) at (5,1.0) {执行与状态回传（闭环更新）};
	\draw[->, thick, draw=black!75] (q1) -- (q2);
	\draw[->, thick, draw=black!75] (q2) -- node[left]{\scriptsize 是} (q3);
	\draw[->, thick, draw=black!75] (q2) -- node[right]{\scriptsize 否} (q4);
	\draw[->, thick, draw=black!75] (q3) -- (q5);
	\draw[->, thick, draw=black!75] (q4) -- (q6);
	\draw[->, thick, draw=black!75] (q5) -- (q7);
	\draw[->, thick, draw=black!75] (q6) -- (q7);
\end{tikzpicture}%
}
\caption{A* 与 D3QN 协同切换决策程序框图}
\label{fig:ch6_switch_flow}
\end{figure}

图\ref{fig:ch6_switch_flow}反映了协同切换的核心逻辑：系统并非在全局与局部策略间静态二选一，而是在统一风险判据下动态分配控制主导权，使规划器在复杂地形中保持可达性与响应性并重。

\subsection{在线重规划鲁棒性与时延容错分析}
\label{subsec:ch6_replan_robustness}

在真实任务中，链路时延与状态噪声会共同影响重规划触发与执行稳定性。设通信时延为 $\Delta t_k$，状态观测噪声为 $\mathbf{n}_k$，则补偿后的局部目标写为
\begin{equation}
\hat{\mathbf{g}}_{k}=\mathbf{g}_{k}+\mathbf{J}_{g}\mathbf{v}_{k}\Delta t_k-\mathbf{K}\mathbf{n}_k,
\label{eq:ch6_delay_compensation_model}
\end{equation}
其中 $\mathbf{J}_{g}$ 为目标点对速度的敏感矩阵，$\mathbf{K}$ 为噪声抑制增益。为评估鲁棒性，引入综合稳定指标
\begin{equation}
\mathcal{R}_{stab}=1-\left(\omega_1\bar{e}_{track}+\omega_2\bar{e}_{yaw}+\omega_3P_{col}\right),
\label{eq:ch6_stability_metric}
\end{equation}
其中 $\bar{e}_{track}$ 为平均轨迹误差，$\bar{e}_{yaw}$ 为平均航向误差，$P_{col}$ 为碰撞概率。

\begin{table}[H]
\centering
\small
\caption{不同扰动条件下在线重规划鲁棒性评估}
\label{tab:ch6_robustness}
\begin{tabular}{p{0.22\textwidth}p{0.16\textwidth}p{0.14\textwidth}p{0.14\textwidth}p{0.14\textwidth}}
\toprule
扰动场景 & 平均轨迹误差(m) & 碰撞概率 & 重规划成功率 & $\mathcal{R}_{stab}$ \\
\midrule
无扰动基准 & 0.42 & 0.010 & 1.000 & 0.936 \\
中等时延扰动 & 0.58 & 0.018 & 0.992 & 0.907 \\
强时延扰动 & 0.81 & 0.031 & 0.978 & 0.861 \\
中等噪声扰动 & 0.64 & 0.022 & 0.988 & 0.892 \\
时延与噪声耦合 & 0.93 & 0.039 & 0.971 & 0.842 \\
\bottomrule
\end{tabular}
\end{table}

从表\ref{tab:ch6_robustness}可见，即使在耦合扰动场景下，重规划成功率仍保持在 $97\%$ 以上，说明补偿机制能够有效控制误差传播。需要指出的是，强扰动条件下轨迹误差与碰撞概率均有上升，提示系统仍需在极端通信退化场景中引入更保守的安全边界。

\begin{algorithm}[H]
\caption{在线增量重规划与容错执行伪代码}
\label{alg:ch6_replan}
\begin{algorithmic}[1]
\REQUIRE 当前状态 $\mathbf{x}_k$，全局路径 $\mathcal{P}_g$，扰动估计 $(\Delta t_k,\mathbf{n}_k)$
\ENSURE 控制动作 $\mathbf{u}_k$
\STATE 计算补偿目标 $\hat{\mathbf{g}}_k$（式\eqref{eq:ch6_delay_compensation_model}）
\STATE 根据式\eqref{eq:ch6_switch_rule}评估策略切换条件
\IF{$\mathcal{S}_k>\tau_{sw}$}
    \STATE 由 D3QN 输出局部动作候选
    \STATE 执行动态安全边界检查
\ELSE
    \STATE 采用 A* 引导动作并执行平滑校正
\ENDIF
\IF{轨迹误差超过阈值或局部碰撞风险升高}
    \STATE 触发增量重规划并更新局部参考段
\ENDIF
\STATE 输出 $\mathbf{u}_k$ 并记录诊断信息
\end{algorithmic}
\end{algorithm}

算法\ref{alg:ch6_replan}强调“补偿-切换-校核-重规划”的时序闭环。该机制将故障检测与决策执行统一到同一周期内，减少了传统串行架构中“先失败再修复”的响应迟滞。

\subsection{多任务长航程综合评估与工程可行性}
\label{subsec:ch6_long_horizon}

为评估算法在长期任务中的稳定收益，本节构建多任务长航程实验，包含“单目标快速到达”“多目标巡检”“动态风险穿越”三类任务。定义综合任务效用为
\begin{equation}
\mathcal{U}_{mission}=w_l\frac{L_{ref}}{L}+w_s(1-P_{col})+w_t\frac{T_{ref}}{T}+w_e\frac{E_{ref}}{E},
\label{eq:ch6_mission_utility}
\end{equation}
其中 $L$、$T$、$E$ 分别表示路径长度、任务时长与能耗，$P_{col}$ 表示碰撞概率。

\begin{table}[H]
\centering
\small
\caption{多任务长航程场景下的综合性能对比}
\label{tab:ch6_long_horizon}
\begin{tabular}{p{0.20\textwidth}p{0.14\textwidth}p{0.14\textwidth}p{0.14\textwidth}p{0.14\textwidth}}
\toprule
任务类型 & 路径长度(m) & 任务时长(s) & 碰撞概率 & $\mathcal{U}_{mission}$ \\
\midrule
单目标快速到达 & 652.7 & 684.2 & 0.010 & 0.921 \\
多目标巡检任务 & 1014.5 & 1098.3 & 0.014 & 0.903 \\
动态风险穿越 & 735.8 & 812.7 & 0.021 & 0.884 \\
多目标+动态风险耦合 & 1126.9 & 1215.4 & 0.028 & 0.861 \\
\bottomrule
\end{tabular}
\end{table}

结果表明，随着任务复杂度提升，综合效用呈可解释下降趋势，但整体仍保持在较高水平。特别是在“多目标+动态风险耦合”场景中，算法未出现不可恢复失效，说明所提混合规划框架具备面向长航程任务的工程可行性。

\begin{figure}[H]
\centering
\resizebox{0.70\textwidth}{!}{%
\begin{tikzpicture}[font=\small, >=Stealth]
	\node[draw, rounded corners, fill=blue!10, align=center, minimum width=7.1cm, minimum height=1.0cm] (m1) at (5,8.6) {任务队列生成与优先级排序};
	\node[draw, rounded corners, fill=cyan!12, align=center, minimum width=7.1cm, minimum height=1.0cm] (m2) at (5,6.8) {全局可达性分析与子目标分配};
	\node[draw, rounded corners, fill=green!12, align=center, minimum width=7.1cm, minimum height=1.0cm] (m3) at (5,5.0) {A-D3QN 混合规划执行};
	\node[draw, rounded corners, fill=orange!14, align=center, minimum width=7.1cm, minimum height=1.0cm] (m4) at (5,3.2) {安全核验与在线重规划触发};
	\node[draw, rounded corners, fill=purple!10, align=center, minimum width=7.1cm, minimum height=1.0cm] (m5) at (5,1.4) {任务完成评估与策略参数回写};
	\draw[->, thick, draw=black!75] (m1) -- (m2);
	\draw[->, thick, draw=black!75] (m2) -- (m3);
	\draw[->, thick, draw=black!75] (m3) -- (m4);
	\draw[->, thick, draw=black!75] (m4) -- (m5);
	\draw[->, thick, draw=blue!70!black] (m5.west) .. controls (1.2,1.2) and (1.2,7.8) .. node[left, align=center]{阶段复盘\\与调参} (m2.west);
\end{tikzpicture}%
}
\caption{多任务长航程规划执行程序框图}
\label{fig:ch6_long_horizon_flow}
\end{figure}

图\ref{fig:ch6_long_horizon_flow}体现了长航程任务中“任务管理-规划执行-安全核验-策略回写”的闭环组织方式。相比单任务场景，该流程增加了任务级调度与阶段复盘环节，使规划器能够在长时序运行中保持稳定性能。

\section{本章小结}
\label{sec:ch6_summary}

本章围绕“数字孪生约束下的月面巡视器路径规划”展开研究，完成了从问题建模、混合算法设计到仿真/地面验证的完整闭环。
主要贡献如下：
\begin{myitemize}
\item 构建了环境建模、状态感知、决策优化、虚实校核一体化规划框架，实现了规划链路与数字孪生系统接口的工程对齐；
\item 提出 A* 与 D3QN 协同的混合规划方法，并通过改进 A* 平滑策略、D3QN 轻量推理与策略切换机制实现全局可达性与局部响应性的兼顾；
\item 通过多复杂度仿真和地面平台验证，定量证明了方法在路径平稳性、碰撞抑制与综合效率方面的优势。
\end{myitemize}

与现有单一全局或单一强化学习方法相比，本章方法的技术优势在于：
将“全局先验约束”与“局部学习决策”统一到同一闭环框架中，避免了纯学习方法可达性不足和纯搜索方法动态适应性不足的问题。

研究局限在于：
当前地面可视化端帧率仍低于工程目标，且部分时延补偿结论基于控制仿真模型而非全实物在环平台。
后续工作将面向三方面展开：
\begin{myitemize}
\item 引入增量式全局重规划与局部策略并行执行，降低复杂场景下秒级规划开销；
\item 开展 D3QN 模型压缩与推理加速，提升端侧部署效率；
\item 构建更高保真地面在环试验链路，系统评估通讯时延、感知误差与执行偏差的耦合影响。
\end{myitemize}
